\section{Liouville and gradient estimates}
\label{sec:hessian-liouville}\index{Liouville theorem!complex Hessian cone}

The gradient estimate comes from rescaling a sequence at points of large
gradient. The background then disappears, leaving a maximal function for
a constant Euclidean cone. We first establish the Liouville theorem that
rules out such a nonconstant limit; its proof is that of
\cref{thm:stw-cone-liouville} for another cone.

For this Euclidean argument allow $1\leq m\leq n$. Write
$H(v)=(v_{i\bar j})$ for the complex Hessian. A continuous real function
$v$ on a domain in $\mathbb C^n$ is admissible for $C_{n,m}$ if
\[
 \operatorname{tr}(PH(v))\geq0
 \quad\text{in distributions for every }P\in C_{n,m}^*,
\]
where $C_{n,m}^*=\{P:\operatorname{tr}(PA)\geq0\text{ for all }A\in C_{n,m}\}$
is the dual cone of Hermitian matrices.
For smooth functions this says $H(v)\in C_{n,m}$. A \emph{lower test}
for $v$ at $z_0$ is a smooth function $\psi$ near $z_0$ with
$\psi\leq v$ near $z_0$ and $\psi(z_0)=v(z_0)$. Call $v$
\emph{maximal} if $H(\psi)(z_0)$ lies outside the interior of $C_{n,m}$
for every lower test $\psi$ at every point $z_0$. Admissibility implies subharmonicity and is
preserved by nonnegative convolution. The cone polynomial $S_m$ increases
under nonnegative matrix addition in the interior.

\begin{lemma}\label{hessian:hessian-cone-stability}
Local uniform limits of admissible maximal functions are admissible and maximal. Suppose $v$ is admissible and maximal on a bounded open set $D$, and $h$ is smooth and admissible there. If both are continuous on $\overline D$ and $h\leq v$ on $\partial D$, then $h\leq v$ on $D$.
\end{lemma}
\begin{proof}
The proof of \cref{lem:cone-maximal-stability-comparison} applies with $C_{n,m}$ in place of the cone $\{P_n(H)\geq0\}$. It uses only that distributional admissibility passes to local uniform limits, that the interior of the cone is open, and that $H+\delta I$ lies in the interior for $H$ in the cone and $\delta>0$.
\end{proof}

\begin{theorem}\label{hessian:hessian-cone-liouville}
Every bounded globally Lipschitz admissible function on $\mathbb C^n$ that is maximal for $C_{n,m}$ is constant.
\end{theorem}
\begin{proof}
For $m=1$, the dual cone consists of the nonnegative multiples of $I$. Distributional subharmonicity and the lower-test condition then say that the function is a viscosity solution of the Laplace equation. The elementary equivalence with harmonicity and bounded harmonic Liouville give the conclusion. For $m=n$, admissibility is distributional plurisubharmonicity. The continuous representative is plurisubharmonic, and bounded plurisubharmonic Liouville applies.
For a fixed $m\geq2$, we argue by induction on $n\geq m$, the case $n=m$ being the plurisubharmonic one.

For a Hermitian $(n-1)$-by-$(n-1)$ matrix $A$ and $m<n$,
\[
 S_j(A\oplus0)=S_j(A),\qquad 1\leq j\leq m.
\]
Consequently $A\oplus0$ belongs to $C_{n,m}$, respectively its interior, exactly when $A$ belongs to $C_{n-1,m}$, respectively its interior. For the closed cones, the nonnegative inequalities characterize the closure: the expansion of $S_j(\lambda+\delta\mathbf1)$ has nonnegative coefficients and a strictly positive $\delta^j$ term. If an admissible function is independent of $z_n$, convolving it and using the block identity shows that the function in the remaining variables is admissible. Extending its lower tests independently of $z_n$ shows that maximality also passes to dimension $n-1$.

With this reduction, the proof of \cref{thm:stw-cone-liouville} applies word for word, with $C_{n,m}$ in place of the cone $\{P_n(H)\geq0\}$, the induction starting at $n=m$, and \cref{hessian:hessian-cone-stability} in place of \cref{lem:cone-maximal-stability-comparison}. Apart from these two inputs, that proof uses only properties shared by $C_{n,m}$: the cone is closed, convex and stable under adding nonnegative matrices; admissibility and maximality are preserved by translations, unitary changes of coordinates and $v\mapsto av+b$ with $a>0$; and admissibility implies subharmonicity and is preserved by nonnegative convolution. For instance, for $v\geq0$ the convolution at a fixed radius of $2v_\epsilon H(v_\epsilon)$ lies in $C_{n,m}$ by convexity.
\end{proof}

A trace bound quadratic in the gradient now bounds the gradient. A
bound linear in the gradient would suffice by
\cref{lem:fixed-trace-gradient-interpolation}, since forms in
$\Gamma_m(\omega)$ have positive trace; the quadratic bound keeps its
form under rescaling, and \cref{hessian:hessian-cone-liouville} excludes
the rescaled limits.

\begin{proposition}\label{hessian:hessian-trace-to-gradient}
Let $\theta$ be a smooth real $(1,1)$-form on $X$, let $2\leq m\leq n$, and let $M,A,B>0$.
There is a uniform bound for $\sup_X|\partial u|_\omega$ over all smooth real $u$ such that $\theta_u=\theta+\ddc u\in\Gamma_m(\omega)$ and
\[
 |u|\leq M,\qquad
 \frac{\theta_u^m\wedge\omega^{n-m}}{\omega^n}\leq B,
 \qquad
 \operatorname{tr}_\omega\theta_u
       \leq A\bigl(1+\sup_X|\partial u|_\omega^2\bigr).
\]
\end{proposition}
\begin{proof}
In local holomorphic coordinates write $\omega=\mathrm i g_{i\bar j}\,\mathrm dz_i\wedge\mathrm d\bar z_j$ and $\theta_u=\mathrm i(b_{i\bar j}+\kappa u_{i\bar j})\,\mathrm dz_i\wedge\mathrm d\bar z_j$, where $\kappa=1/(2\pi)$ comes from the normalization of $\ddc$.
Since $m\geq2$, every eigenvalue of a form in $\Gamma_m(\omega)$ has modulus at most its trace, by the proof of \cref{hessian:weighted-cone-trace}. Thus the assumed trace estimate bounds every eigenvalue of $\theta_u$ relative to $\omega$ in absolute value by $A(1+\sup|\partial u|^2)$.

Suppose, contrary to the claim, that a sequence $u_j$ has maxima $R_j=\sup|\partial u_j|_\omega\to\infty$, attained at $x_j$. Passing to a subsequence, let $x_j\to x_\infty$, and work in a fixed chart, changing coordinates linearly so that $g(x_\infty)=I$. If $y_j$ is the coordinate of $x_j$, set
\begin{align*}
 v_j(z)&=u_j(y_j+z/R_j),\qquad g_j(z)=g(y_j+z/R_j),\\
 E_j(z)&=R_j^{-2}b(y_j+z/R_j)+\kappa H(v_j)(z).
\end{align*}
These are defined on balls with radii tending to infinity. On each fixed ball $g_j\to I$, the smooth background term tends to zero, and $E_j$ is uniformly bounded by the eigenvalue estimate. In particular $H(v_j)$ and its Euclidean trace are uniformly bounded.
Interior estimates for the fixed Euclidean Laplacian give local uniform $W^{2,q}$ bounds for every finite $q$. A diagonal subsequence converges in $C^1$ on compact subsets to a bounded globally Lipschitz $v$, with $|\partial v(0)|=1$.

The Hermitian matrices $g_j^{-1/2}E_jg_j^{-1/2}$ belong to the interior of $C_{n,m}$ and differ from $\kappa H(v_j)$ by a term tending uniformly to zero on fixed balls. Pairing with any fixed matrix in the dual cone and passing in distributions proves that $v$ is admissible.
By homogeneity of $S_m$, the rescaled Hessian densities satisfy
\[
 \frac{S_m(g_j^{-1/2}E_jg_j^{-1/2})}{\binom nm}
   =R_j^{-2m}\frac{\theta_{u_j}^m\wedge\omega^{n-m}}{\omega^n}
      (y_j+z/R_j)\leq BR_j^{-2m}.
\]
We prove maximality with lower tests. If a smooth lower test $\psi$ at $z_0$ had $H(\psi)(z_0)$ in the interior, choose a small $\delta>0$ so that $H(\psi)-\delta I$ remains there near $z_0$. Strictifying the contact by subtracting $\delta|z-z_0|^2$ gives interior minimum points $z_j\to z_0$ for $v_j$ minus this test. At those points
\[
 E_j\geq R_j^{-2}b+\kappa(H(\psi)-\delta I).
\]
The matrix on the right, measured with $g_j$, is uniformly inside the cone. Monotonicity of $S_m$ under nonnegative addition gives a fixed positive lower bound for the left side of the preceding equation. This contradicts $BR_j^{-2m}\to0$.
Hence $v$ is maximal. \Cref{hessian:hessian-cone-liouville} makes it constant, contradicting $|\partial v(0)|=1$.
\end{proof}
